\section{TFPT / Universalraum: Operationssatz, nativer Grundzustand und
Feldwörterbuch}\label{tfpt-universalraum-operationssatz-nativer-grundzustand-und-feldwuxf6rterbuch}

\textbf{Forschungsfortsetzung v1.6.3 · 15. September 2026}

Diese Runde bearbeitet die vier Prioritäten aus v1.6.2 §14 in der dort
festgelegten Reihenfolge: den tatsächlich verfügbaren Operationssatz
exakt abgrenzen, den berichteten nativen Grundzustand mit seinem
ursprünglichen Hamilton- und Sektorvertrag reproduzieren und darauf die
geladene Antwort berechnen, das relativistische Feldwörterbuch am
tatsächlichen Tensor festhalten. Eine räumliche Skalierungsrechnung ist
ausdrücklich \textbf{nicht} Gegenstand; sie trägt erst, wenn Zustand,
Operationen und Feldtyp stehen. Haupt- und Update-PDF v1.6 bleiben
unverändert. Kein vollständiges T1-T8-Tor wird geschlossen.

\subsection{1. Ergebnis in einem Absatz}\label{ergebnis-in-einem-absatz}

Der N=64-Sektor des nativen Modells ist exakt als Loch-Paarerzeugung auf
dem gefüllten Zustand realisiert; seine Krylov-Kette ist bis zu 15 252
960 Einträgen exakt berechnet, ihre Normen ν₁\ldots ν₄ sind auf zwei
völlig unabhängigen Wegen (Besetzungsbild und Onishi/Wick-Spurnetzwerke)
identisch. Die Kette schließt bei zwei Bosonen \textbf{exakt nicht}
(Defekt 2,9·10⁻⁵), das Modell ist also nicht allein durch seine Normen
lösbar, numerisch aber fast. Rigorose Ritz-Obergrenzen, Bosonbesetzung,
Überlapp, Verschränkungsuntergrenze und die vollständige
Ein-Teilchen-Antwort (Z\_h + Z\_add = 1) liegen vor. Die
Worker-Behauptung eines eindeutigen N=64-Grundzustands bis g/Δ = 1/20
ist mit Casimir-Schranken nur bis g/Δ ≤ 0,0251 beweisbar und bei 1/20
\textbf{nicht reproduziert} (nicht widerlegt). Die Kommutanten-Leiter
zeigt exakt, wie viel unsichtbare Freiheit jede Operationsstufe
beseitigt: zwei Kontrollen lassen 1,44·10⁹ Freiheitsgrade in N=3, die
volle innere Symmetrie sieben. Das Feldwörterbuch ist am Tensor
entschieden: der Vermittler kann nur Vektor oder antisymmetrischer
Tensor sein; ein einheitliches Vektor-Wörterbuch existiert nicht
(Trägergraph nicht bipartit, 5-Zyklus {[}16,45,0,30,37{]}); der einzige
einheitliche Typ ist (1,0)⊕\allowbreak{}(0,1), ein B\_\{μν\}, kein Spin 2.

\subsection{2. Unveränderte Quelle und
Konventionen}\label{unveruxe4nderte-quelle-und-konventionen}

\[\adjustbox{max width=\linewidth}{$\displaystyle 
H=\Delta N_b+g\sum_{A=1}^{60}(b_A^\dagger P_A+P_A^\dagger b_A),\qquad
P_A=\sum_{i<j}W_{A,ij}f_jf_i,\qquad N=N_f+2N_b,\qquad WW^{\mathsf T}=8I_{60}.
$}\]

Der Tensor wird repo-lokal gelesen (SHA-256
\codeword{3f00a0892157a691e4ef26920c3702772c7bf76a250fad98ff04a079bb64b763}),
der Clock-Konstruktor ist gepinnt (\codeword{9bf99de7…}).
\codeword{common.py} prüft: 60 Kanäle × 8 disjunkte Paare, jede Mode in
15 Kanälen; Gewichtserhaltung auf allen 480 Paaren; die
Casimir-Identität 8WᵀW + 𝒞\_S + 𝒞\_C − 120 I = 0 auf allen 2016
Zweiteilchenzuständen (𝒞 = −ΣL², d.~h. viermal der Standard-Casimir);
induzierte Bosongeneratoren X\_B = WΛ²(X)Wᵀ/8 gaußganzzahlig und exakt
kovariant; Clock-Lift Periode sechs, Slot-Permutation (2,0,1,4,3),
Vorzeichen −1, WΛ²G\_F = G\_B W.

\subsection{3. Operationssatz und
Kommutant}\label{operationssatz-und-kommutant}

Gemessen wird die Dimension des Kommutanten der von einer
Operationsstufe erzeugten Algebra auf dem Sektor: kleiner Kommutant
bedeutet weniger unsichtbare Freiheit. Das Spektralzertifikat für N=3
ist reproduziert (S = C₃C₃ᵀ wird von ∏(S−r), r ∈ \{0,7,10,12\}, exakt
annulliert; Multiplizitäten 64, 2880, 576, 320 aus rationalen
Projektorspuren).

\subsubsection{3.1 Exakte Irreduzibel-Zerlegung
(Höchstgewichtsvektoren)}\label{exakte-irreduzibel-zerlegung-huxf6chstgewichtsvektoren}

{\def\LTcaptype{none} % do not increment counter
\begin{longtable}[]{@{}
  >{\raggedright\arraybackslash}p{(\linewidth - 8\tabcolsep) * \real{0.1875}}
  >{\raggedright\arraybackslash}p{(\linewidth - 8\tabcolsep) * \real{0.1875}}
  >{\raggedleft\arraybackslash}p{(\linewidth - 8\tabcolsep) * \real{0.2500}}
  >{\raggedright\arraybackslash}p{(\linewidth - 8\tabcolsep) * \real{0.1875}}
  >{\raggedright\arraybackslash}p{(\linewidth - 8\tabcolsep) * \real{0.1875}}@{}}
\toprule\noalign{}
\begin{minipage}[b]{\linewidth}\raggedright
Sektor
\end{minipage} & \begin{minipage}[b]{\linewidth}\raggedright
Block
\end{minipage} & \begin{minipage}[b]{\linewidth}\raggedleft
Dim
\end{minipage} & \begin{minipage}[b]{\linewidth}\raggedright
Zerlegung unter so(10)⊕\allowbreak{}su(4)
\end{minipage} & \begin{minipage}[b]{\linewidth}\raggedright
(𝒞\_S, 𝒞\_C)
\end{minipage} \\
\midrule\noalign{}
\endhead
\bottomrule\noalign{}
\endlastfoot
N=2 & helle Paare ↔ Bosonen & 60 & (10, 6) & (36, 20) \\
N=2 & dunkle Paare & 1956 & (120, 10) ⊕\allowbreak{} (126, 6) & (84,36), (100,20) \\
N=3 & dunkle Tripel = ker C₃ & 37888 & (560, 20″) ⊕\allowbreak{} (1200, 20) ⊕\allowbreak{} (672,
4̄) & (117,63), (141,39), (165,15) \\
N=3 & ker S & 64 & (16, 4̄) & (45, 15) \\
N=3 & λ = 12 & 320 & (16, 20) & (45, 39) \\
N=3 & λ = 10 & 576 & (144, 4̄) & (85, 15) \\
N=3 & λ = 7 & 2880 & (144, 20) & (85, 39) \\
\end{longtable}
}

Alle Multiplizitäten sind eins; Σ m·dim stimmt in jedem Block mit der
Weyl-Dimensionsformel. Höchstgewichte in verdoppelten Koordinaten, z. B.
(560,20″) ↔ ({[}3,3,1,1,−1{]},{[}3,3,3{]}), (672,4̄) ↔
({[}3,3,3,3,3{]},{[}1,1,−1{]}).

\subsubsection{3.2 Kommutanten-Leiter}\label{kommutanten-leiter}

{\def\LTcaptype{none} % do not increment counter
\begin{longtable}[]{@{}
  >{\raggedright\arraybackslash}p{(\linewidth - 4\tabcolsep) * \real{0.2727}}
  >{\raggedleft\arraybackslash}p{(\linewidth - 4\tabcolsep) * \real{0.3636}}
  >{\raggedleft\arraybackslash}p{(\linewidth - 4\tabcolsep) * \real{0.3636}}@{}}
\toprule\noalign{}
\begin{minipage}[b]{\linewidth}\raggedright
Verfügbare Operationen
\end{minipage} & \begin{minipage}[b]{\linewidth}\raggedleft
N=2
\end{minipage} & \begin{minipage}[b]{\linewidth}\raggedleft
N=3
\end{minipage} \\
\midrule\noalign{}
\endhead
\bottomrule\noalign{}
\endlastfoot
A: \{X, N\_b\} (die zwei bisher geprüften Kontrollen) & 3 829 536 & 1
444 233 216 \\
B: A + Clock (Ordnung 6) & 641 760 & 240 742 144 \\
Cc: A + 15 su(4)-Generatoren & 30 376 & 2 247 168 \\
Cs: A + 45 so(10)-Generatoren & 172 & 1 648 \\
C: A + alle 60 Generatoren & 3 & 7 \\
D: B + C & 3 & 7 \\
B(H\_N) & 1 & 1 \\
\end{longtable}
}

Für A ist die Blockstruktur M₃₇₈₈₈ ⊕\allowbreak{} M₆₄ ⊕\allowbreak{} (I₂⊗M₂₈₈₀) ⊕\allowbreak{} (I₂⊗M₅₇₆) ⊕\allowbreak{}
(I₂⊗M₃₂₀); für B folgen die Zahlen aus exakten Spuren der Clock-Potenzen
je Block (z. B. Multiplizitäten 6384/6280/6280/6384/6280/6280 auf dem
Dunkelraum); für die Lie-Stufen aus Σ m\_i². Der Clock ist auf
\textbf{keinem} irreduziblen Block ein Skalar (er verschiebt Gewichte;
Zeugen im JSON), also ein äußeres Element der Symmetrie, nicht in ihrer
zusammenhängenden Gruppe. D = C gilt trotzdem, weil alle Multiplizitäten
eins sind und der Clock jeden Block in sich abbildet.

\textbf{Folgerung.} Mehrfachwiederholung der zwei Kontrollen beseitigt
nichts; der Clock allein reduziert um einen Faktor sechs; erst die
innere Symmetrie als Operationssatz kollabiert die Freiheit auf je einen
Skalar pro (Energieblock, Irrep). so(10) ist dabei die wirksamere
Hälfte. Ob der Compiler diese Symmetrieoperationen tatsächlich
bereitstellt, ist damit nicht bewiesen; ihr Effekt ist exakt beziffert.

\subsection{4. Der native
N=64-Grundzustand}\label{der-native-n64-grundzustand}

\subsubsection{4.1 Lochbild}\label{lochbild}

Mit h\_i = f\_i† ist \textbar F⟩ = f₀†\ldots f₆₃†\textbar0⟩ das
Lochvakuum und als Operatoridentität P\_A = −Q\_A† mit Q\_A† =
Σ\_\{i\textless j\}W\_\{A,ij\}h\_i†h\_j†. Unter (−1)\^{}\{N\_b\} ist H
äquivalent zu

\[\adjustbox{max width=\linewidth}{$\displaystyle 
H_{\rm hole}=\Delta N_b+g\sum_A(b_A^\dagger Q_A^\dagger+Q_Ab_A),\qquad N_h=2N_b.
$}\]

Die Teilchen-Loch-Unitäre trägt das Vorzeichen s(μ) = (−1)\^{}\{i+j\} je
Zweilochzustand; die naive Komplement-Abbildung ohne s(μ) ist keine
Identität (Negativkontrolle bestanden).

\subsubsection{4.2 Exakte Krylov-Kette}\label{exakte-krylov-kette}

v\_n = (T†)ⁿ\textbar F⟩ mit T† = Σ\_A b\_A†Q\_A†:

{\def\LTcaptype{none} % do not increment counter
\begin{longtable}[]{@{}rrrl@{}}
\toprule\noalign{}
n & Einträge in v\_n & ν\_n = ‖v\_n‖² & ν\_n/ν\_\{n−1\} \\
\midrule\noalign{}
\endhead
\bottomrule\noalign{}
\endlastfoot
0 & 1 & 1 & - \\
1 & 480 & 480 & 480 \\
2 & 108 240 & 439 680 & 916 = 4·229 \\
3 & 15 252 960 & 575 078 400 & 299520/229 \\
4 & (nur Spurnetzwerk) & 952 296 652 800 & ≈ 1656,0 \\
5 & (nur Spurnetzwerk) & siehe \codeword{norms\_by\_traces.json} & \\
\end{longtable}
}

ν₁\ldots ν₃ sind auf zwei unabhängigen Wegen identisch (Besetzungsbild
mit exakten Ganzzahlamplituden; Spurnetzwerke, §6). Exakte Guards: T v₁
= 480 v₀; T v₂ = 916 v₁ (das Ein-Boson-Singulett (60⊗Λ²64̄)\^{}G ist
eindimensional); ⟨v₂\textbar T v₃⟩ = ν₃ (Adjungiertheit); v₁, v₂ sind
Singuletts unter je drei so(10)- und su(4)-Generatoren; ⟨v₁\textbar Σ\_A
Q\_AQ\_A†\textbar v₁⟩ = 458·480, was die Formel Σ\_A Q\_AQ\_A† = 480 −
(15N\_h + C\_std\^{}\{Loch\})/2 mit dem Loch-Casimir 14 der (10,6)
bestätigt - der Lochanteil von v₁ ist kein Singulett, nur der
Gesamtzustand.

\subsubsection{4.3 Schließung}\label{schlieuxdfung}

T v₃ = α v₂ + w₂ mit α = 299520/229, w₂ ⊥ v₂ und exakt

\[\adjustbox{max width=\linewidth}{$\displaystyle 
\|w_2\|^2=\|Tv_3\|^2-\frac{\nu_3^2}{\nu_2}=\frac{5001523200}{229}\approx 2{,}18\cdot10^7,
\qquad \frac{\|w_2\|^2}{\|Tv_3\|^2}=2{,}90\cdot10^{-5}.
$}\]

Der Zwei-Boson-Singulettraum ist also mehr als eindimensional; die Kette
ist \textbf{exakt nicht geschlossen}, und das Modell ist nicht allein
durch die ν\_n lösbar. Numerisch ist die Korrektur vernachlässigbar: w₂
koppelt nur an v₃ mit g·√(‖w₂‖²/ν₃) = 0,195 g.

\subsubsection{4.4 Ritz-Obergrenzen}\label{ritz-obergrenzen}

In der orthonormierten Basis \{v₀\ldots v\_K, w₂\} ist H (Δ = 1)
tridiagonal plus w₂-Zweig: Diagonale n bzw. 2; Nebendiagonalen g·4√30,
g·2√229, g·48√29770/229, \ldots; w₂-Kopplung g·3√598377/11908.

{\def\LTcaptype{none} % do not increment counter
\begin{longtable}[]{@{}lllll@{}}
\toprule\noalign{}
g/Δ & K=1 & K=2 & K=3 (+w₂) & −480g² \\
\midrule\noalign{}
\endhead
\bottomrule\noalign{}
\endlastfoot
1/20 & −0,704159 & −0,985178 & −1,094232 & −1,2 \\
1/40 & −0,241620 & −0,288824 & −0,295357 & −0,3 \\
1/100 & −0,045894 & −0,047851 & −0,047894 & −0,048 \\
1/400 & −0,002991 & −0,003000 & −0,003000 & −0,003 \\
\end{longtable}
}

Bei g/Δ = 1/20 ist die Kette bei K=3 nicht konvergiert (die
Stufenkopplungen liegen bei 1,1Δ\ldots1,6Δ); die Verlängerung mit ν₄, ν₅
steht in \codeword{native\_ground\_state.json} (Schlüssel
\codeword{ritz}). Bei g/Δ ≤ 1/40 ist die Konvergenz praktisch erreicht.

\subsubsection{4.5 Observablen auf dem Ritz-Zustand (g/Δ = 1/20,
K=3)}\label{observablen-auf-dem-ritz-zustand-gux3b4-120-k3}

⟨N\_b⟩ = 0,8421 (Hellmann-Feynman ∂E/∂Δ = 0,8421);
\textbar⟨F\textbar Ω⟩\textbar² = 0,397; Bosonzahlverteilung 0,397 /
0,397 / 0,172 / 0,034; Shannon-Untergrenze der Boson-Loch-Verschränkung
1,66 bit (rigoros, weil die reduzierte Bosondichte in N\_b blockdiagonal
ist). Die Worker-Angaben „Überlapp \textgreater{} 1/4`` und
„Verschränkung \textgreater{} 0,811 bit`` sind damit konsistent, aber
auf einem nichtkonvergierten Zustand.

\subsubsection{4.6 Geladene
Ein-Teilchen-Antwort}\label{geladene-ein-teilchen-antwort}

Für jede Mode r (Guard r ∈ \{0, 5, 63\}; Transitivität der
Quellsymmetrie):

\[\adjustbox{max width=\linewidth}{$\displaystyle 
Z_h=\|f_r\Omega\|^2=1-\frac{\langle N_b\rangle}{32}=0{,}97368,\qquad
Z_{\rm add}=\|f_r^\dagger\Omega\|^2=\frac{\langle N_b\rangle}{32}=0{,}02632,\qquad
Z_h+Z_{\rm add}=1.
$}\]

Mittlere Anregungsenergien: Entnahme ε\_h = 0,0311 Δ, Addition ε\_add =
1,1497 Δ. Kompositnorm Z\_χ = 23⟨N\_b⟩/480 = 0,0404 (die N=4-Referenz
aus v1.6.2 ergab 0,0301; verschiedene Zustände).

\subsubsection{4.7 Sektorauswahl und die
Worker-Behauptung}\label{sektorauswahl-und-die-worker-behauptung}

Untergrenzen aus der Casimir-Identität und quadratischer Ergänzung, H ≥
(1−θ)ΔN\_b − (15g²/2θΔ)N\_f, ausgewertet als max\_θ min\_\{N\_b\} je
Sektor:

{\def\LTcaptype{none} % do not increment counter
\begin{longtable}[]{@{}llll@{}}
\toprule\noalign{}
g/Δ & Ritz E(64) & schärfster Konkurrent & N=64 bewiesen global? \\
\midrule\noalign{}
\endhead
\bottomrule\noalign{}
\endlastfoot
1/400 & −0,0029996 & N=63: −0,0029531 & ja \\
1/100 & −0,047894 & N=63: −0,04725 & ja \\
1/40 & −0,295357 & N=63: −0,2953125 & ja \\
1/20 & −1,094232 & N=66: −1,2; N=62: −1,1625 & \textbf{nein} \\
\end{longtable}
}

Bewiesener Bereich mit diesen Schranken: g/Δ ≤ 0,0251. Die Behauptung
„eindeutiger Spin(10)×SU(4)-Singulett-Grundzustand für 0 \textless{}
\textbar g\textbar/Δ ≤ 1/20`` ist damit \textbf{nicht reproduziert und
nicht widerlegt}; bei 1/20 müsste E(64) \textless{} −1,2 gegen N ≥ 66
und \textless{} −1,1625 gegen N = 62 gezeigt werden. μ* = −E/64 = 0,0171
Δ liegt unter Δ/50; die frühere Aussage „μ = Δ/50 wählt den leeren
Zustand`` bleibt konsistent.

\subsection{5. Feldwörterbuch am tatsächlichen
Tensor}\label{feldwuxf6rterbuch-am-tatsuxe4chlichen-tensor}

\subsubsection{5.1 Symmetrieregel}\label{symmetrieregel}

Für Grassmann-Felder ψ\_\{I,a\} ist Σ\_\{IJ\}Σ\_\{ab\}
W\_\{IJ\}K\_\{ab\}ψ\_\{I,a\}ψ\_\{J,b\} = ½Σ(T − Tᵀ) mit T\_\{(Ia),(Jb)\}
= W\_\{IJ\}K\_\{ab\}. Weil W antisymmetrisch ist, verschwindet die
Kopplung genau dann, wenn K antisymmetrisch ist; \textbf{nur der
symmetrische Teil des Spinorkerns koppelt}. Auf allen 60 Zeilen explizit
geprüft:

{\def\LTcaptype{none} % do not increment counter
\begin{longtable}[]{@{}
  >{\raggedright\arraybackslash}p{(\linewidth - 4\tabcolsep) * \real{0.3333}}
  >{\raggedright\arraybackslash}p{(\linewidth - 4\tabcolsep) * \real{0.3333}}
  >{\raggedright\arraybackslash}p{(\linewidth - 4\tabcolsep) * \real{0.3333}}@{}}
\toprule\noalign{}
\begin{minipage}[b]{\linewidth}\raggedright
Zuordnung
\end{minipage} & \begin{minipage}[b]{\linewidth}\raggedright
Kern K
\end{minipage} & \begin{minipage}[b]{\linewidth}\raggedright
Kopplung
\end{minipage} \\
\midrule\noalign{}
\endhead
\bottomrule\noalign{}
\endlastfoot
gleichhändige Weyl, Boson Skalar (0,0) & ε antisym. & 0 \\
gleichhändige Weyl, Boson (1,0) & σ\^{}\{μν\}ε sym. & ≠ 0 (Rang 32 je
Komponente) \\
Dirac/Majorana: Skalar C, Pseudoskalar Cγ⁵, Axialvektor Cγ\^{}μγ⁵ &
antisym. & 0 \\
Dirac/Majorana: Vektor Cγ\^{}μ, Tensor Cσ\^{}\{μν\} & sym. & ≠ 0 \\
\end{longtable}
}

Der Vermittler ist damit Vektor oder antisymmetrischer Tensor, niemals
Skalar, Pseudoskalar oder Axialvektor. Er trägt U(1)\_N-Ladung 2; ein
masseloser geladener Vektor ist kein konsistentes freies Eichfeld
(Vorbehalt, kein Satz).

\subsubsection{5.2
Chiralitätsgradierungen}\label{chiralituxe4tsgradierungen}

Der Trägergraph (64 Moden, 480 Kanten, Grad 15, dreiecksfrei) ist
\textbf{nicht bipartit}; ein ungerader Zyklus ist {[}16, 45, 0, 30,
37{]}. Es gibt daher keine globale Chiralitätsaufteilung, unter der alle
60 Bosonen Vektoren wären. Für jede der 44 G-kovarianten Gradierungen
ist jeder Kanal rein (nie gemischt), aber der Bosonmultiplett zerfällt:

{\def\LTcaptype{none} % do not increment counter
\begin{longtable}[]{@{}
  >{\raggedright\arraybackslash}p{(\linewidth - 8\tabcolsep) * \real{0.1667}}
  >{\raggedleft\arraybackslash}p{(\linewidth - 8\tabcolsep) * \real{0.2222}}
  >{\raggedleft\arraybackslash}p{(\linewidth - 8\tabcolsep) * \real{0.2222}}
  >{\raggedright\arraybackslash}p{(\linewidth - 8\tabcolsep) * \real{0.1667}}
  >{\raggedleft\arraybackslash}p{(\linewidth - 8\tabcolsep) * \real{0.2222}}@{}}
\toprule\noalign{}
\begin{minipage}[b]{\linewidth}\raggedright
Gradierung
\end{minipage} & \begin{minipage}[b]{\linewidth}\raggedleft
Vektor-Kanäle
\end{minipage} & \begin{minipage}[b]{\linewidth}\raggedleft
Tensor-Kanäle
\end{minipage} & \begin{minipage}[b]{\linewidth}\raggedright
Stabilisator
\end{minipage} & \begin{minipage}[b]{\linewidth}\raggedleft
Kommutant der 64
\end{minipage} \\
\midrule\noalign{}
\endhead
\bottomrule\noalign{}
\endlastfoot
Identität & 0 & 60 & 60 & 1 \\
Pati-Salam-Spinorsplit (10 Varianten) & 24 & 36 & 36 = so(6)⊕\allowbreak{}so(4)⊕\allowbreak{}su(4)
& 2 \\
Farbsplit (3) & 40 & 20 & 52 & 2 \\
gemischt (30) & 32 & 28 & 28 & 4 \\
\end{longtable}
}

Ein einheitlicher Lorentztyp für alle 60 Vermittler erzwingt die
gleichhändige Zuordnung und damit (1,0)⊕\allowbreak{}(0,1): ein antisymmetrischer
Tensor B\_\{μν\}, kein symmetrisches Spin-2-Feld. Kinetische Terme:
einer pro irreduziblem Block des Stabilisators, keine
Familienaufspaltung 4 → 1+3. Ein Modenindex ist kein Raumpunkt.

